\chapter{绪\hspace{6pt}论}
\label{chap:intro}

\section{研究工作的背景与意义}

\subsection{海洋立体观测对水空直接通信的需求}

海洋观测、海底资源勘探、海上应急搜救和近海安防正在从单一平台作业转向由水下航行器、海面浮标、无人艇和空中平台共同组成的立体网络。网络中的水下节点负责近距离感知和数据采集，空中节点负责大范围中继与回传，两者之间需要一条能够直接穿越海面的信息通道。目前，水下节点与外界通信普遍采用“水声通信—海面浮标—无线电/卫星”的接力方式。水声链路带宽通常在千比特每秒量级，传播时延为秒级，海面中继节点的布放与维护也增加了系统的复杂度和暴露面。

蓝绿波段光在海水中存在低损耗窗口，水下无线光通信（Underwater Wireless Optical Communication，UWOC）能够在数米至数十米距离上提供兆比特每秒至吉比特每秒的速率\citing{kaushal2016uowc,zeng2017survey,arnon2010uowc}。若光束能够从水下直接穿过海面到达空中接收机，水下节点就可以不经过海面中继而与空中平台交换数据。已有实验在静水和波动水面条件下验证了水空光链路的可行性\citing{chen2017airwater,sun2019wavy}，说明这一方向具备工程基础。水下航行器在执行任务时需要上传高清图像、声呐数据和状态信息，水空直接光链路能够缩短数据回传的时延，并降低对海面中继设施的依赖。

\subsection{水空直接光链路面临的物理约束}

水空直接光链路的上行方向从高折射率的海水进入低折射率的空气，这一方向决定了其物理约束与空水下行链路不同，主要体现在以下三个方面。

第一，界面全反射限制了可用的入射角范围。海水折射率约为 1.33，空气折射率约为 1，水—空气界面的临界角约为 $48.75^\circ$。当水下发射机与海面中继点之间的连线偏离竖直方向较大时，光束轴向入射角超过临界角，弹道光几乎全部被反射回水中。即使入射角略小于临界角，Snell 定律也会把空气侧出射角放大到接近掠射方向：水中 $40^\circ$ 入射对应空气中约 $58.8^\circ$ 出射。空气侧接收机必须布置在远离海面中继点水平位置的低仰角方向上，这对空中平台的悬停位置和接收视场提出了较强约束。

第二，斜入射带来强像散和光斑展宽。折射角对入射角的导数在临界角附近迅速增大，光束在入射面内的发散被界面放大，空气侧光斑沿一个方向被拉长。在本文的仿真场景中，水下 10\,m 斜程、$40^\circ$ 入射的高斯光束穿过平静水面后，在 10\,m 高度接收平面上的二阶矩光斑长短轴分别为 93.0\,cm 和 46.9\,cm，轴比接近 2。光斑展宽直接降低了有限孔径接收机的捕获功率。

第三，收发两端都存在随机失准。水下航行器受海流和推进控制误差影响，其位置和航向缓慢漂移\citing{weng2022pointing}；海面上的任何漂浮器件都随波浪升沉、横摇和纵摇；空中接收平台存在位置抖动和指向角误差。接收孔径相对光斑的位置偏移与接收视场的角度误差共同决定瞬时捕获功率，并与海洋湍流、大气湍流引起的光强闪烁叠加，形成误码率平台和中断事件。

\subsection{液晶智能超表面在水空链路中的作用}

智能超表面（Reconfigurable Intelligent Surface，RIS）通过对入射波逐单元施加可编程的相位或幅度调制来重塑传播信道，已在射频通信中得到广泛研究\citing{basar2019ris,wu2019irs}，并逐步扩展到自由空间光通信\citing{najafi2021irsfso}。光频段的相位梯度超表面服从广义 Snell 定律\citing{yu2011generalized}：界面切向的相位梯度直接改变出射光的切向波矢，从而在不改变介质折射率的条件下改变折射方向。

液晶相位调制器是一类成熟的透射式光学相控器件\citing{mcmanamon1996opa,mcmanamon2009review}。对液晶层施加电压可以改变液晶分子的取向和等效折射率，通过微电极阵列加载阶梯状或闪耀状相位分布即可实现光束偏转。Kim 等研制的大面积液晶光束偏转器在 14.4\,mm $\times$ 14.4\,mm 有效区内实现了 7200 个独立控制通道\citing{kim2020lcdeflector}，为构建米级面板的拼接式液晶 RIS 提供了器件参考。

将透射式液晶 RIS 漂浮在海面上、下层玻璃贴近水面，可以在光束进入空气之前施加横向相位梯度。这一结构为水空直接链路带来三项能力：（1）抵消一部分水中入射光的切向波矢，使空气侧出射方向向竖直方向偏转，降低对空中接收平台仰角的要求；（2）在光束到达玻璃—空气界面之前减小其切向波矢，使原本超过裸水面临界角的光线仍能透射进入空气；（3）通过逐单元相位分配控制接收平面上的光斑尺寸和形状，在接收功率与抖动容忍度之间进行权衡。

围绕这一结构建立可计算、可验证的系统模型，需要依次回答以下问题：海面漂浮的 RIS 在风浪中如何运动，运动如何改变光束在 RIS 下表面的入射几何；水下光束经过吸收和散射后在 RIS 下表面形成怎样的光强分布与相位；受器件偏转角、角度量化和单元填充率约束的实际相位能在接收平面上形成怎样的光场；接收机的位置抖动、角度误差、有限视场和湍流衰落如何共同决定误码率和中断概率；不同优化目标下的相位方案之间存在怎样的性能权衡。本文围绕上述问题展开研究，为水空直接无线光通信系统的设计与评估提供建模工具和数值依据。

\section{本文的主要研究内容与贡献}

本文以“水下发射机—海面漂浮液晶 RIS—空中接收机”构成的水空直接无线光链路为研究对象，按照“坐标与器件—平台运动与相位设计—入射与接收端复光场—动态海面下的链路性能”的顺序建立系统模型，并在风浪驱动的粗糙海面上评估 RIS 对链路性能的增强效果。主要研究内容与贡献如下。

（1）建立了统一的系统坐标体系和贴水安装的液晶 RIS 器件方案。定义了世界坐标系、水平参考系、RIS 本体与面板坐标系、发射机光轴坐标系和接收孔径坐标系及其变换链；在均匀水体、高斯光源和吸收型 RIS 边界下，把辐射输运方程化为积分方程，给出以单次散射反照率为公比收敛的碰撞级数解；把分层介质中的 Fermat 路径、面积 Jacobian 与 Fresnel 透射组合为具有振幅和相位的复光场传递关系。以 Neumann 方向谱生成随机海面并审计离散方差。参考大面积液晶光束偏转器，给出了 $1\,\mathrm{m}\times1\,\mathrm{m}$、$40\times40$ 单元玻璃夹层液晶 RIS 的单元排布、电极周期、偏转角预算与角度量化，论证了下层玻璃贴水安装对规避裸水面全反射的作用。

（2）建立了海面漂浮 RIS 的六自由度运动模型与受器件约束的相位设计方法，提出了预测调相与鲁棒光斑结合的调相方式。完成了环形浮框的质量预算与静水配平，建立了以单位四元数描述姿态的六自由度刚体方程，将薄板辐射/绕射载荷、侧壁压力及其解析偏航力矩、辐射记忆状态空间模型纳入同一时域积分框架。在给定位姿的面板坐标中，由全部单元的公共可达方向确定接收机位置，以入射能量的白化径向分位数为各单元分配平顶、高斯或径向 Weibull 型鲁棒目标落点，经分层 Fermat 路径反解、限幅量化和常相位校准得到可加载相位。在此基础上，以运动模型预测相位生效时刻的位姿并据此设计相位，把位姿预测残差折算为接收平面上的等效抖动，再按等效抖动选择鲁棒光斑宽度。

（3）建立并验证了 RIS 入射光场与 RIS 调制后接收端复光场的计算模型。将 RIS 下表面的光强分解为弹道分量与散射分量：弹道分量在有限束腰源和任意倾斜面板几何下给出了精确的面积 Jacobian 表达式，散射分量以完整 Henyey--Greenstein 相函数下的碰撞次数展开定义，采用强制首次散射的无偏估计器求值，在 18 个冻结快照上与独立蒙特卡洛结果的差异处于后者统计噪声水平。对任意可加载相位，以二阶光程展开下的逐单元衍射积分、Chirp-$z$ 变换、Parseval 功率守恒搬运和像素内相干平均计算接收端弹道光强，叠加散射光的非相干贡献，并给出大接收距离下的频谱窗口与邻近单元相干处理；在 10--50\,m、3--12\,m/s 的 228 个相位上，复光场的接收功率与独立蒙特卡洛结果的差异在 10\% 以内。

（4）在冻结的动态粗糙海面下系统评估了 RIS 对水空链路的性能增强。同一组海浪谐波同时给出 RIS 位姿和直达光穿过的海面，无 RIS 直达链路以粗糙海面蒙特卡洛计算，Steering-only、Flat top、Robust ABER/OP 和 Max mean power 四个 RIS 方案在 RIS 可达范围内共用同一个抖动的接收机。结果表明：直达链路在接收距离 30\,m 及以上时 90\,dBm 以内达不到 $10^{-3}$ 的平均误码率，最优 RIS 方案只需约 42\,dBm，且所需功率几乎不随接收距离和风速变化；Robust ABER/OP 与 Flat top 是明显优于其他方案的两个方法，在 30\,m 处比 Steering-only 低约 6.5\,dB，比最大平均功率方案低约 14\,dB，其中鲁棒光斑在强海洋湍流下的中断概率上进一步领先。静水机理分析表明，最优光斑半径约为位置抖动标准差的 2.4--2.9 倍，入射角超过临界角后 RIS 是唯一可用的上行通道。

\section{本论文的结构安排}

本文共分为六章，各章内容安排如下。

第一章为绪论。介绍水空直接无线光通信的研究背景，分析上行链路在界面全反射、斜入射像散和收发失准方面的物理约束，说明液晶 RIS 在水空链路中的作用，并概述本文的研究内容与贡献。

第二章为水空直接光链路的基础理论与液晶 RIS 器件设计。建立全文统一的坐标系，介绍海水固有光学特性、辐射输运方程的碰撞级数解、海浪谱与随机海面生成、分层介质与粗糙海面上的折射及复光场传递、液晶相位调制与广义 Snell 定律，给出液晶 RIS 的器件设计，并介绍湍流衰落、指向误差和通信性能指标。

第三章研究海面漂浮 RIS 的六自由度运动建模与相位设计。完成浮框配平，建立风浪激励下的六自由度刚体运动模型并给出运动结果；在给定位姿下建立受器件约束的有限宽度光斑相位设计方法与基准方案，提出预测调相与鲁棒光斑结合的调相方式。

第四章研究 RIS 入射光场与调制后接收端复光场的建模与验证。建立发射端运动与跟踪模型，推导弹道分量与完整 HG 散射分量的光束扩展函数并与独立蒙特卡洛结果比较；在高斯光束与标量衍射理论的基础上给出液晶层入射复光场，并对任意相位建立 RIS 调制后接收平面的复光场表达式及其蒙特卡洛验证。

第五章研究动态海面下 RIS 辅助水空链路的性能增强。建立无 RIS 直达链路的粗糙海面蒙特卡洛模型和计及抖动、视场与湍流的接收性能模型，在多种接收距离、风速和湍流强度下比较五条链路的平均误码率与中断概率，选出最优相位方法，并在静水条件下分析光斑宽度、接收端扰动和入射角的作用机理。

第六章为全文总结，归纳本文的主要研究成果。
